% Modernised reading — fonds Alexandre Grothendieck, shelfmark 102 (pages 1-13).
%
% This is an INTERPRETATION, not a transcription. It re-reads the pages in
% today's notation and vocabulary, states the hypotheses the manuscript leaves
% implicit, and drops the critical apparatus so that the mathematics reads
% straight through. Everything the brackets and underlines carried moves into
% a footnote. It is held to being mathematically correct as it stands; where
% the manuscript is loose or mistaken the true statement appears and the
% footnote says what the page has. In any dispute of substance the
% transcription governs: whoever cites Grothendieck cites batch-01.fr.tex.
%
% Scope: a single document for the whole folder. The shelfmark holds one
% batch (pages 1-13), read whole before any of this was written. Page 1 is an
% orange cover carrying « Recollements »; pages 7 (blank) and 9 (a printed
% verso with no hand on it) are not transcribed.
%
% Spine. Five pieces around one programme, the « étude générale des
% foncteurs (Sch)° -> (Ens) » of the inventory's title: pages 2-6, a
% typescript « Pour EGA VI » on gluing objects of a site along a class M of
% morphisms (schemes and Artin's algebraic spaces as the two examples);
% page 8, four questions for Artin on descent; page 10, a plan in eight items
% for the study of properties of morphisms of functors; page 11, a fragment
% of a proof through the Hilbert scheme; pages 12-13, a criterion for a
% morphism of functors to be an open immersion, and a sheet of keywords. The
% links between the pieces (the « variétés » of page 8 as the objects of
% pages 2-6; condition (i) of page 12 as the « (A) » of page 13 and the
% marginal of page 8) are THIS READING'S.
%
% Notation fixed for the whole folder. S is the site of pages 2-6 (and, on
% pages 8 and 12, a base scheme: kept apart by context, said once in the
% spine section); S~ its topos of sheaves, M the class of arrows, M~ its
% extension, S~_M the glued objects, S(X,M) the small site of X. « fpqc »,
% « lfp » (localement de présentation finie). A morphism of functors is said
% to have a property P stable under base change when every base change to a
% scheme has P (principle of page 10).
%
% Substantive departures from the pages, with pages.
%  - Page 2, d): read as locality on the source for M-coverings (the
%    parenthesis « i.e. par a) et c) » makes the pr_1 condition equivalent to
%    u_i in M). Stated so.
%  - Page 6: open immersions do NOT satisfy d) (X u X -> X is a
%    counterexample), so the first example works only in the variant the
%    margin of page 2 proposes (no d), saturate in M~); S~_M is then still
%    the category of schemes. Ours, footnoted.
%  - Pages 2-3: the reduction of case 2 to case 1 is known only through the
%    transcription's summary; a subsheaf of a representable sheaf is not
%    representable in general, and the footnote says what the reduction needs.
%  - Page 4: « le plus petit sous-ensemble de Ob S » read Ob S~ (as the
%    transcription already notes).
%  - Page 11: « plat de prés. finie sur Y » read over Y'.
%  - Page 12: the proof sketched for the open-immersion criterion is ours;
%    the page states it without proof.
%
% Announced and never established in the folder: the existence half of the
% Proposition of page 2 (« laissé au lecteur »); (v) of page 4 (« au
% rédacteur »); the Proposition of page 6 (no proof); answers to the four
% questions of page 8; all eight items of the plan of page 10 except the
% criterion of page 12; the theorem whose proof page 11 continues (its
% statement is not in the folder).
%
% Group. The folder is filed among « Géométrie et topologie combinatoire »
% (69-102), but nothing in it meets the combinatorial folders 81, 88, 89. It
% meets folder 91, whose page 4 checks the same conditions (fpqc sheaf,
% commutation with filtered colimits and with adic limits of artinian rings)
% for a representability theorem; cross-referenced by its pages.
%
% Literature named in footnotes (EGA plan, SGA 3 XI, SGA 4 VIII, Artin,
% Knutson, the Stacks Project) is cited from memory, not checked against the
% sources. This reading was made on Opus 5.5 and has not been measured
% against the reference reading of folder 115 (made on Opus 5).
%
% Pass: Opus 5.5 (claude-opus-5-5), 2026-10-10 — first pass, unchecked against the pages by a human.
\documentclass[11pt,a4paper]{article}
\input{../preamble/grothendieck}

\folder{102}
\pages{1}{13}
\dating{[après 1967]}
\watermark{Édition de démonstration\\interprétation personnelle de l'œuvre}
\foldertitle{Etude générale des foncteurs (Sch)° → (Ens) : tapuscrit annoté (s.d.), notes manuscrites (s.d.) — lecture modernisée du dossier entier}

\begin{document}

\section*{Résumé}

\begin{resume}
Une carte routière d'un pays se fait en recollant des cartes de régions :
chacune est simple, et toute l'information est dans la façon dont elles se
recouvrent. La géométrie algébrique de Grothendieck fabrique ses objets de la
même manière. Les morceaux simples sont les spectres d'anneaux — des objets
définis par un anneau commutatif — et un schéma est ce qu'on obtient en les
recollant le long de parties ouvertes. Ces treize pages posent la question
dans sa généralité : qu'obtient-on si l'on recolle les mêmes morceaux le long
d'une autre classe de flèches que les inclusions d'ouverts ?

La réponse du tapuscrit des pages 2 à 6, prévu pour un chapitre d'EGA qui n'a
jamais paru, tient en un procédé et deux exemples. On se place dans la
catégorie de \emph{tous} les « espaces généralisés » — les faisceaux sur un
site, c'est-à-dire des foncteurs qui à chaque morceau simple associent
l'ensemble de ses applications vers l'espace —, on y prolonge la classe de
flèches choisie, et l'on retient les objets qui se laissent couvrir par des
morceaux simples à l'aide de flèches de cette classe. Avec les immersions
ouvertes on retrouve les schémas ; avec les morphismes étales — des flèches
qui sont des « revêtements locaux », comme un revêtement en topologie mais
sans exiger qu'il soit trivial au-dessus d'un ouvert — on trouve les
\emph{espaces algébriques} que Michael Artin venait d'introduire, objets plus
souples que les schémas et dans lesquels certaines constructions, impossibles
pour les schémas, deviennent possibles.

Le reste du dossier est manuscrit et rapide. Quatre questions pour Artin
demandent quand une « donnée de descente » — la recette pour fabriquer un
objet sur une base à partir d'un objet sur une base plus grande — est
effective. Un plan en huit points dresse la liste des propriétés des
morphismes de schémas (être une immersion, être lisse, étale, propre, séparé…)
qu'il faudrait étendre aux foncteurs, avec un principe pour le faire : une
propriété doit être locale, et coïncider avec la définition reçue quand le
foncteur est un schéma. Une page de démonstration, dont le théorème manque,
utilise le schéma de Hilbert pour rendre réduites les fibres d'une famille.
Enfin un critère : un sous-foncteur qui se comporte bien vis-à-vis des
limites et des épaississements infinitésimaux est une immersion ouverte.

Ce qui se dégage de la manière : ne jamais construire un objet à la main
quand on peut le définir par ce qu'il représente. Le schéma, l'espace
algébrique, l'immersion ouverte sont ici des propriétés de foncteurs, que
l'on vérifie une à une comme on coche une liste — et la feuille de mots
encadrés de la page 13 en est littéralement la liste.

\keywords{gluing, subcanonical site, local property of morphisms, algebraic space, étale morphism, small étale site, fpqc descent, effective descent datum, Artin approximation, representable morphism, valuative criterion, formally étale morphism, limit-preserving functor, Hilbert scheme, geometrically reduced fibre, open immersion criterion, Artin's representability criteria}
\end{resume}

\pagerange{2}{13}
\section*{Le fil du dossier, et les conventions}

La chemise (page 1) porte le seul mot « Recollements ». Les pages 7 et 9 ne
portent rien de sa main. Le reste se répartit en cinq morceaux :

\begin{itemize}
  \item \textbf{Pages 2 à 6 : tapuscrit « Pour EGA VI ».} Recollement d'objets
    d'un site $S$ le long d'une classe $M$ de flèches ; prolongement $\widetilde{M}$
    de $M$ au topos ; la catégorie $\widetilde{S}_M$ des objets recollés, ses
    propriétés de stabilité ; schémas et espaces algébriques ; le topos
    localisé au-dessus d'un objet. Annoté à la main.
  \item \textbf{Page 8 : « Questions pour Artin ».} Effectivité de la descente
    pour des schémas propres sur une base hensélienne, descente fpqc des
    « variétés » et des petits faisceaux étales.
  \item \textbf{Page 10 : un plan.} Huit rubriques de propriétés des morphismes
    de foncteurs, et le principe de vérification.
  \item \textbf{Page 11 : fragment d'une démonstration.} Le théorème démontré
    n'est pas dans le dossier.
  \item \textbf{Pages 12 et 13 : immersions ouvertes.} Un critère, puis une
    feuille de mots-clés encadrés.
\end{itemize}

Le titre de l'inventaire, « Étude générale des foncteurs
$(\mathrm{Sch})^{\circ} \to (\mathrm{Ens})$ », est celui de la page 10 et
couvre l'ensemble : les objets recollés des pages 2 à 6 sont des foncteurs de
ce type, et les pages suivantes cherchent comment parler de leurs morphismes.
Le lien entre le tapuscrit et la page 8 est de cette lecture : les
« variétés au sens généralisé, sans exiger que la diagonale soit une
immersion » de la page 8 sont les objets que le tapuscrit obtient avec les
morphismes étales.

\textbf{Conventions, valables pour tout le dossier.} Dans les pages 2 à 6, $S$
désigne un site ; aux pages 8 et 12, la même lettre désigne un schéma de base,
comme sur les pages, et le contexte suffit à les distinguer. $\widetilde{S}$
est le topos des faisceaux d'ensembles sur le site $S$ ; « fpqc » est la
topologie fidèlement plate et quasi-compacte ; « lfp » abrège « localement de
présentation finie ». Un foncteur $(\mathrm{Sch})^{\circ} \to (\mathrm{Ens})$
— un préfaisceau sur la catégorie des schémas, contravariant comme tout
préfaisceau — est identifié au schéma qu'il représente quand il en représente
un ; un morphisme de foncteurs $f : G \to F$ est dit avoir une propriété $P$
\emph{stable par changement de base} si, pour tout schéma $T$ et tout
$T \to F$, le morphisme $G \times_F T \to T$ est un morphisme de schémas ayant $P$
— c'est le principe de la page 10.

\textbf{Le dossier dans son groupe.} Il est classé parmi les dossiers 69 à 102,
« Géométrie et topologie combinatoire », mais n'a rien de commun avec les cartes,
hexagones et découpages des dossiers 81, 88 et 89, et sa datation propre,
« après 1967 », le place avant la période du groupe. Il rencontre en revanche
le dossier 91 : la page 4 de celui-ci vérifie, pour un théorème de
représentabilité, les conditions mêmes qu'énonce la page 12 ci-dessous
(faisceau fpqc, commutation aux limites inductives filtrantes d'anneaux et aux
limites projectives adiques d'anneaux artiniens).

\pagerange{2}{6}
\section*{Recollements d'objets d'un site (pages 2 à 6)}

Le tapuscrit est titré « Pour EGA VI »\footnote{Le plan des \emph{Éléments de
géométrie algébrique} annoncé en 1960 prévoyait, sauf erreur de mémoire, un
chapitre VI consacré à la descente et aux méthodes générales de construction
des schémas ; il n'a jamais paru. Le rattachement du tapuscrit à ce chapitre
est la page ; le contenu prévu du chapitre est cité de mémoire.}. Comme pour
les tapuscrits d'autres dossiers, la transcription donne les énoncés et les
interventions de sa main, et résume les démonstrations tapées ; la présente
lecture ne peut donc rien affirmer de ces démonstrations au-delà de ce résumé.

\pagerange{2}{2}
\subsection*{Le site et la classe de flèches}

Soit $S$ un site dont la topologie est moins fine que la topologie canonique
— un site \emph{sous-canonique}, ce qu'il appelle « standard »\footnote{La
précision « standard (i.e.\ à topologie moins fine que la canonique) » est une
insertion interlinéaire. Elle garantit que les foncteurs représentables sont
des faisceaux, de sorte que le plongement de Yoneda identifie $S$ à une
sous-catégorie pleine de $\widetilde{S}$ ; tout ce qui suit en dépend.}. Soit
$M$ une classe de flèches de $S$ telle que :

\begin{itemize}
  \item[a)] les flèches de $M$ sont quarrables (leurs produits fibrés avec
    toute flèche de $S$ existent dans $S$), et $M$ est stable par changement de
    base ;
  \item[b)] $M$ contient les identités et est stable par composition ;
  \item[c)] $M$ est \emph{locale sur le but} : si $u : X \to Y$ est une flèche
    de $S$ et $(Y_j \to Y)$ une famille couvrante telle que chaque
    $X \times_Y Y_j \to Y_j$ soit dans $M$, alors $u$ est dans $M$ ;
  \item[d)] $M$ est \emph{locale sur la source pour les recouvrements par
    $M$} : si $(u_i : X_i \to X)$ est une famille couvrante de flèches de $M$
    et si les composés $u u_i : X_i \to Y$ sont dans $M$, alors $u : X \to Y$ est
    dans $M$.
\end{itemize}

La forme donnée à d) est une reformulation\footnote{La condition d) est
entièrement réécrite à la main sur la page 2. La version tapée demandait que les
$X_i \to Y$ soient dans $M$ pour des « $Y$-morphismes » $u_i$ ; la version
retenue demande les $u u_i$ dans $M$ et les projections
$X_i \times_X X_j \to X_i$ dans $M$, cette dernière clause étant introduite par
« i.e.\ par a) et c) ». Nous lisons la parenthèse ainsi : les projections
$X_i \times_X X_j \to X_i$ sont les changements de base de $u_j$ le long de la
famille couvrante $(u_i)$, de sorte que, par a) et c), demander qu'elles soient
dans $M$ revient à demander $u_j \in M$. Le « $S$ » qui suit « morphismes dans »
est surchargé d'un trait et a pu être changé.}. C'est ce qu'on appelle
aujourd'hui une propriété \emph{locale sur la source} : elle vaut pour les
morphismes étales dans la topologie étale, et c'est elle qui permettra plus loin
de dire « étale » d'un morphisme d'espaces algébriques sans ambiguïté.

En marge, il envisage de s'en passer : ne pas imposer d) à $M$, mais saturer
plus tard dans le prolongement $\widetilde{M}$, en décrétant qu'une flèche est
dans $\widetilde{M}$ si et seulement si elle devient dans $M$ après composition
avec un recouvrement par des flèches de $M$\footnote{Note écrite de biais dans la
marge gauche de la page 2, reliée à d) par une accolade, en regard aussi du
début de la proposition suivante, dont elle propose de reformuler la condition
$\tilde{e}$).}. On verra à la page 6 que cette variante est exactement ce que
demande son premier exemple.

\pagerange{2}{3}
\subsection*{Le prolongement de $M$ au topos}

\textbf{Proposition.} Il existe une unique classe $\widetilde{M}$ de flèches du
topos $\widetilde{S}$ qui satisfait aux analogues $\tilde{a}$) à $\tilde{d}$)
de a) à d), et de plus :
\begin{itemize}
  \item[$\tilde{e}$)] une flèche de $S$ est dans $\widetilde{M}$ si et
    seulement si elle est dans $M$ ;
  \item[$\tilde{f}$)] si $u : F \to X$ est dans $\widetilde{M}$ avec $X$
    représentable, il existe une famille couvrante $(u_i : X_i \to F)$ de
    flèches de $\widetilde{M}$ avec les $X_i$ représentables\footnote{Le
    « $S$ » de « $X \in S$ » est surchargé à la main et se lit mal. À la fin de
    $\tilde{f}$), quelques mots tapés sont biffés, où l'on devine
    « $uu_i$ … $\widetilde{M}$ ». Les étiquettes $\tilde{a}$ à $\tilde{f}$ sont
    remarquées d'un tilde à la main, sur les pages 2 et 3.}.
\end{itemize}

Le tapuscrit démontre l'unicité en montrant que l'appartenance d'une flèche
$u : F \to G$ à $\widetilde{M}$ est déterminée par les conditions. Par
$\tilde{c}$), on se ramène au cas où $G = X$ est représentable. Si la diagonale
$F \to F \times_X F$ est représentable, $u$ est dans $\widetilde{M}$ si et
seulement s'il existe une famille couvrante $(u_i : X_i \to F)$, $X_i \in S$,
telle que les $u u_i$ et les projections $X_i \times_F X_j \to X_i$ soient dans
$M$ ; les $X_i \times_F X_j$ sont représentables comme images inverses de la
diagonale de $F$ dans les $X_i \times_X X_j$, de sorte que la condition a un
sens. Le cas général est ramené au précédent\footnote{D'après le résumé de la
transcription, le cas 2º) observe que $Z_{ij} = X_i \times_F X_j$ est un
sous-faisceau du faisceau représentable $X_i \times_X X_j$, image inverse de la
diagonale. Un sous-faisceau d'un faisceau représentable n'est pas représentable
en général ; la réduction demande donc que l'on sache déjà, par récurrence, que
$Z_{ij} \to X_i$ est dans $\widetilde{M}$, ce qui, par $\tilde{f}$), le couvre
par des objets de $S$. Dans le cas étale, c'est le point où la théorie des
espaces algébriques « s'amorce » : la représentabilité de la diagonale se
démontre après coup au lieu d'être supposée. Faute du texte tapé, nous ne
pouvons dire comment le tapuscrit règle ce point.}. Le tapuscrit laisse au
lecteur la vérification que la classe ainsi décrite satisfait bien à
$\tilde{a}$) à $\tilde{f}$) : l'existence n'est pas démontrée dans le
dossier.

\pagerange{3}{3}
\subsection*{Les objets recollés}

\textbf{Définition.} Un faisceau $F \in \widetilde{S}$ \emph{s'obtient par
recollement d'objets de $S$ à l'aide de morphismes dans $M$} s'il existe une
famille couvrante $(u_i : X_i \to F)$ avec $X_i \in S$ et
$u_i \in \widetilde{M}$\footnote{Les mots « par recollement d'objets de $S$ à
l'aide de morphismes dans $M$ » sont soulignés à la main.}. On note
$\widetilde{S}_M$ la sous-catégorie pleine de $\widetilde{S}$ formée de ces
faisceaux.

Un tel $F$ est le quotient de $\coprod_i X_i$ par la relation d'équivalence
$R = \coprod_{i,j} X_i \times_F X_j$, dont les deux projections sont dans
$\widetilde{M}$ : c'est la forme la plus générale d'une « carte par
recollement », les $X_i$ étant les cartes et $R$ les changements de cartes.

\pagerange{4}{5}
\subsection*{Stabilité}

\textbf{Proposition.}\footnote{Une première rédaction, « Proposition 1) »,
commencée au bas de la page 3 et poursuivie par deux lignes au haut de la
page 5 (verso de la page 4), est biffée ; elle portait sur une famille couvrante
$F_i \to F$ à projections dans $\widetilde{M}$, et la conclusion « il faut et il
suffit que pour tout $i$ … » s'interrompt. La version retenue est celle de la
page 4.}
\begin{itemize}
  \item[(i)] Tout objet de $S$ est dans $\widetilde{S}_M$.
  \item[(ii)] Si $u : F \to G$ est dans $\widetilde{M}$ et $G \in \widetilde{S}_M$,
    alors $F \in \widetilde{S}_M$. La réciproque vaut si $u$ est couvrant
    (épimorphisme de faisceaux).
  \item[(iii)] Si $u : F \to G$ est couvrant et si
    $\mathrm{pr}_1 : F \times_G F \to F$ est dans $\widetilde{M}$, alors
    $G \in \widetilde{S}_M$ si et seulement si $F \in \widetilde{S}_M$.
  \item[(iv)] Une somme $\coprod F_i$ est dans $\widetilde{S}_M$ si et seulement
    si tous les $F_i$ le sont.
  \item[(v)] Si $S$ a les produits de deux objets (resp.\ les produits fibrés,
    resp.\ un objet final, resp.\ les limites projectives finies), il en est de
    même de $\widetilde{S}_M$\footnote{Le numéro (iv) surcharge un « (iii) »
    tapé ; « $\mathrm{pr}_1$ » dans (iii) est ajouté à la main ; (v) est retapé
    en interligne sur une première rédaction biffée (« Soit $I$ une catégorie
    finie, telle que … »).}.
\end{itemize}

Pour (ii) : on couvre $G$ par des $Y_j \to G$ de $\widetilde{M}$, on change de
base, on couvre chaque $F \times_G Y_j$ par $\tilde{f}$), et l'on compose par
$\tilde{b}$). Le point (iii) en résulte, $u$ étant dans $\widetilde{M}$ par
$\tilde{c}$) puisque son changement de base le long de lui-même, couvrant,
l'est. Pour (iv), la transcription résume l'argument tapé : si $X \to \coprod F_i$
est dans $\widetilde{M}$, les $X \times_F F_i$ se couvrent par $\tilde{f}$), et
ces morceaux couvrent les $F_i$. Le point (v) est « laissé au rédacteur », avec
une mise en garde : le noyau d'un double flèche ne se comporte sans doute pas
de même, car un morphisme de doublets de $\widetilde{S}$, ou même de $S$, qui
induit des morphismes de $\widetilde{M}$ sur les sommets n'induit pas
nécessairement un morphisme de $\widetilde{M}$ sur les noyaux.

\textbf{Corollaire.} $\mathrm{Ob}\,\widetilde{S}_M$ est la plus petite classe
d'objets de $\widetilde{S}$ qui contient $\mathrm{Ob}\,S$ et est stable par
sommes et par passage au quotient par une relation d'équivalence dont la
première projection est dans $\widetilde{M}$\footnote{Le tapuscrit porte « le plus
petit sous-ensemble de $\mathrm{Ob}\,S$ », sans tilde ; le sens veut
$\widetilde{S}$. « sommes et par » est ajouté à la main.}. Par suite,
$\widetilde{S}_M$ ne change pas si l'on remplace $S$ par une sous-catégorie
$S'$ de $\widetilde{S}$ contenue dans $\widetilde{S}_M$, et $M$ par
$M' = \widetilde{M} \cap \mathrm{Fl}\,S'$\footnote{L'énoncé s'entend avec $S'$ munie de la topologie induite par celle
de $\widetilde{S}$, les objets recollés étant pris dans le même topos
$\widetilde{S}$. Pour les prendre dans le topos des faisceaux sur $S'$, il faut
que celui-ci s'identifie à $\widetilde{S}$ — c'est le cas, par le lemme de
comparaison de SGA 4, dès que $S'$ contient $S$, ce qui est la situation de
l'exemple suivant (schémas affines dans schémas). Précision ajoutée par cette
lecture.}.

Le corollaire résulte de la description d'un objet recollé comme quotient de
$\coprod X_i$ (section précédente), et de (iii) et (iv). C'est l'énoncé qui dit
qu'un recollement de recollements est encore un recollement.

\pagerange{6}{6}
\subsection*{Schémas et espaces algébriques}

\textbf{Exemples.} Soit $S$ la catégorie des schémas affines, munie d'une
topologie sous-canonique plus fine que la topologie de Zariski.

\emph{Immersions ouvertes.} Avec pour $M$ les immersions ouvertes, ou seulement
les immersions $\mathrm{Spec}(A_f) \to \mathrm{Spec}(A)$, la catégorie
$\widetilde{S}_M$ est équivalente à celle des schémas. Il faut pour cela
utiliser la variante que propose la marge de la page 2, sans la condition d) :
les immersions ouvertes ne satisfont pas à d)\footnote{La page dit « $M$
l'ensemble des immersions ouvertes (ou, en l'absence de d), les immersions de
la forme $\mathrm{Spec}(A_f) \to \mathrm{Spec}(A)$) », ce qui suggère que les
immersions ouvertes satisfont à d). Elles ne le font pas : si $Y$ est affine et
$X = Y \sqcup Y$, les deux inclusions $Y \to X$ forment une famille couvrante
d'immersions ouvertes dont les composés avec $u : X \to Y$ sont l'identité, et
$u$ n'est pas une immersion ouverte. Avec d) imposé et la description de
$\widetilde{M}$ de la page 2, la flèche $u$ serait même dans $\widetilde{M}$,
contredisant $\tilde{e}$). La remarque est de cette lecture ; la conclusion
« on retrouve les schémas » reste juste, puisqu'elle ne dépend que de la
définition de $\widetilde{S}_M$ : un faisceau recouvert par des sous-foncteurs
ouverts affines est un schéma.}. Les morphismes de $\widetilde{M}$ sont alors
les immersions ouvertes de faisceaux.

\emph{Morphismes étales.} Avec une topologie plus fine que la topologie étale
(et sous-canonique, par exemple fppf ou fpqc) et pour $M$ les morphismes étales,
qui satisfont à d), on trouve les \emph{espaces algébriques} d'Artin, obtenus
« en recollant des spectres d'anneaux par des morphismes étales »\footnote{Les
espaces algébriques sont introduits par M.~Artin à la fin des années 1960 et
exposés par D.~Knutson (\emph{Algebraic spaces}, 1971), avec une condition de
quasi-compacité sur la diagonale ; la définition sans condition de séparation
est celle qu'on trouve aujourd'hui par exemple dans le Stacks Project. Le
tapuscrit n'impose aucune condition de séparation, et c'est la seconde qu'il
faut entendre. Références citées de mémoire. La page nomme Artin et met
l'expression entre guillemets.}. Par le corollaire, on obtient le même résultat
en partant des schémas au lieu des schémas affines. Les flèches de
$\widetilde{M}$ s'appellent alors morphismes \emph{étales} d'espaces
algébriques, la topologie étant la topologie étale\footnote{Cette dernière
phrase est ajoutée à la main, à l'encre bleue, à la suite de la ligne tapée.}.

\textbf{Proposition.} Supposons que $M$ satisfasse la condition, qu'il appelle
« draconienne » : si $u : X \to Y$ et $v : Y \to Z$ sont des flèches de $S$ avec
$v \in M$ et $vu \in M$, alors $u \in M$. Pour $X \in S$, soit $S(X,M)$ la
sous-catégorie pleine de $S_{/X}$ formée des $X' \to X$ dans $M$ — toutes ses
flèches sont alors dans $M$ —, munie de la topologie induite. Alors la
sous-catégorie pleine de $\widetilde{S}_{/X}$ formée des $F \to X$ qui sont dans
$\widetilde{M}$ est un topos, équivalent, par restriction à $S(X,M)$, au topos
$S(X,M)^{\sim}$. Un tel $F$ est d'ailleurs dans $\widetilde{S}_M$, par (ii) de
la page 4\footnote{La remarque entre crochets « d'où résulte
$F \in \mathrm{Ob}\,\widetilde{S}_M$ ! » est de sa main, à l'encre bleue. Une
première rédaction de la conclusion, rayée et encadrée, énonçait que le composé
$(\widetilde{S}_M)_{/X} \to \widetilde{S}_{/X} \simeq (S_{/X})^{\sim} \to
S(X,M)^{\sim}$ est une équivalence ; elle était fausse telle quelle, puisque
$(\widetilde{S}_M)_{/X}$ contient des $F \to X$ qui ne sont pas dans
$\widetilde{M}$, et la correction restreint justement aux flèches de
$\widetilde{M}$. Le dernier terme y était tapé « $S(M,X)$ ».}.

La condition draconienne est satisfaite par les immersions ouvertes et par les
morphismes étales. Dans le second cas, $S(X,M)$ est le \emph{petit site étale}
de $X$, et la proposition dit que les faisceaux sur le petit site étale sont
exactement les espaces algébriques étales sur $X$ : c'est un énoncé connu
aujourd'hui sous cette forme. Le dossier n'en contient pas de
démonstration\footnote{Aucune démonstration ne suit l'énoncé sur la page 6, qui
est la dernière du tapuscrit.}.

\pagerange{8}{8}
\section*{Questions pour Artin (page 8)}

La page est une liste de questions, sans réponses\footnote{Le titre, souligné,
est de sa main. Le dossier ne contient aucune réponse d'Artin.}. Dans les trois
premières, $S$ est le spectre d'un anneau local noethérien
\emph{hensélien}\footnote{« hens. » est doublement souligné ; « strictement »
et « noeth. » sont biffés devant. La transcription laisse incertaine la lecture
« plat » pour l'hypothèse sur $S' \to S$, précédée d'un mot biffé illisible.} et
$s$ son point fermé.

\begin{itemize}
  \item[1º)] Soit $S' \to S$ surjectif, et sans doute plat, avec $S'$
    noethérien, et $X'$ un $S'$-schéma \emph{propre} muni d'une donnée de
    descente relativement à $S' \to S$\footnote{« propre » remplace « loc.\ de
    t.f. », biffé.}. Si la donnée de descente induite sur la fibre $X'_s$ est
    effective, l'est-elle sur $X'$ ? Et dans le cas $S' = \widehat{S}$, le
    complété ?
  \item[2º)] Peut-on faire la descente fpqc des « variétés », ou des faisceaux
    relativement de présentation finie\footnote{L'insertion se lit « ou
    relation loc.\ de prés.\ finie » ; « relativement » est une autre lecture.},
    par exemple de parties d'un complété d'un anneau local ? Autrement dit, pour
    un faisceau fpqc lfp, la propriété d'être une variété est-elle locale pour
    la topologie fpqc sur $S$ ?
  \item[3º)] Si $X$ est une variété propre sur $S$ dont la fibre fermée $X_s$
    est représentable, $X$ est-il représentable ? Si la réponse à 2º) était
    affirmative, les questions 1º) et 3º) seraient équivalentes.
  \item[4º)] Descente fpqc des petits faisceaux étales — question liée à la
    caractérisation des foncteurs $(\mathrm{Sch})^{\circ}_{/S} \to (\mathrm{Ens})$
    qui proviennent de petits faisceaux étales sur $S$. Une réponse affirmative
    à 2º), pour des variétés au sens généralisé où l'on n'exige pas que la
    diagonale soit une immersion, en entraînerait une à 4º).
\end{itemize}

En marge, deux groupes de questions : chercher des exemples de $X$ propre sur
$S$ local complet avec $X_s$ représentable et $X$ non représentable ; et une
remarque, en partie illisible, liant le problème à la commutation aux limites
projectives adiques d'anneaux artiniens\footnote{Notes écrites de biais dans la
marge gauche, en regard de 2º), puis de 3º) et 4º) ; plusieurs mots en sont
incertains ou illisibles, en particulier « le comm. », « adiques artiniennes »,
« propres » et « repr. ». « 4º » est entouré et relié à la marge.} — la
condition même qui reparaît à la page 12.

Lecture. Les « variétés » sont ici des foncteurs du type que construit le
tapuscrit : la précision « sans exiger que la diagonale soit une immersion »
est ce qui sépare les espaces algébriques des schémas, dont la diagonale est
toujours une immersion. La question 4º) rejoint alors la proposition de la
page 6 : un petit faisceau étale sur $S$ est un espace algébrique étale sur $S$,
de sorte que sa descente est un cas particulier de 2º). Le contexte de 1º) et
3º) est que la descente fpqc des schémas, même propres, n'est pas toujours
effective\footnote{Un exemple célèbre est celui de Hironaka : une variété
propre munie d'une action libre d'un groupe d'ordre $2$ dont le quotient n'est
pas un schéma, ce qui donne une donnée de descente étale non effective. Le
passage d'une base complète à une base hensélienne, demandé à la fin de 1º), est
la question à laquelle répond le théorème d'approximation d'Artin (1969).
Références citées de mémoire ; la page ne cite ni l'un ni l'autre.} ; on
sait aujourd'hui que les faisceaux étales, eux, se descendent le long des
morphismes fidèlement plats et quasi-compacts\footnote{SGA 4, exposé VIII,
cité de mémoire et non vérifié.}.

\pagerange{10}{10}
\section*{Un plan pour l'étude des foncteurs (page 10)}

La page ne porte pas de titre ; c'est un plan numéroté, où une flèche ramène la
rubrique 7 entre 5 et 6. On le donne dans l'ordre que la flèche indique\footnote{La
page garde l'ordre 1, 2, 3, 4, 5, 7, 6, 8 ; la flèche de la marge le corrige.
La frontière entre 5 et 7 est incertaine : la ligne est surchargée, et
« Morphismes », d'une encre plus sombre, est écrit par-dessus.}, en
nommant à côté ce que chaque rubrique est devenue.

\begin{enumerate}
  \item[1.] \emph{Morphismes représentables} et \emph{$P$-représentables}
    (représentables par des morphismes de schémas ayant la propriété $P$) ;
    l'abus de langage qui dit « immersion ouverte, lisse, étale » d'un morphisme
    de foncteurs représentable par de tels morphismes ; le critère de
    représentabilité par recollement de morceaux, avec un renvoi à SGA 3,
    exposé XI, 3.5\footnote{La page écrit « SGAD XI 3.5 ». SGAD est le
    séminaire Demazure–Grothendieck, SGA 3 ; nous n'avons pas vérifié le
    contenu du numéro cité.}.
  \item[2.] Faisceaux et morphismes quasi-compacts, quasi-séparés ; morphismes
    lfp et de présentation finie\footnote{Les crochets de la rubrique 2
    signalent, d'après la marge, de « petits compléments » (lecture incertaine ;
    « conditions » est possible).}. Pour un foncteur, être lfp revient à
    commuter aux limites inductives filtrantes d'anneaux — c'est la
    caractérisation qu'utilise la page 12.
  \item[3.] Morphismes formellement lisses, formellement nets, formellement
    étales, et variantes infinitésimales ; les critères standard.
  \item[4.] Critères pour qu'un morphisme soit un monomorphisme, un
    bimorphisme (à la fois mono et épi), une immersion, une immersion ouverte.
  \item[5.] Conditions de séparation : diagonale représentable par des
    immersions, par des immersions quasi-compactes, par des immersions de
    présentation finie ; diagonale représentable par des immersions fermées ;
    relations avec les critères valuatifs.
  \item[7.] Foncteurs sur une catégorie de schémas au-dessus d'une base, et
    « espaces induits »\footnote{« induits » est une lecture incertaine, et le
    mot qui suit « foncteurs » est biffé et illisible, de même que la base de
    $(\mathrm{Sch})_{/\cdot}$.}. Morphismes surjectifs, morphismes
    radiciels — ceux dont la diagonale est surjective, caractérisation exacte
    pour les morphismes de schémas. Une immersion ouverte surjective est un
    isomorphisme.
  \item[6.] Conditions de propreté.
  \item[8.] Conditions de connexité.
\end{enumerate}

\textbf{Principe des vérifications.} Pour définir une propriété $P$ des
morphismes de faisceaux sur $(\mathrm{Sch})$, deux exigences : a) $P$ doit être
locale ; b) si $f : F \to G$ est représentable, $P$ doit coïncider avec la
définition déjà admise pour les morphismes de schémas, c'est-à-dire : pour tout
schéma $X$ et tout $X \to G$, le morphisme de schémas $X \times_G F \to X$ a la
propriété $P$. C'est la règle qu'on suit aujourd'hui pour étendre aux espaces
algébriques et aux champs les propriétés des morphismes de schémas. Elle n'est
cohérente que pour des propriétés stables par changement de base et locales
sur le but pour la topologie choisie\footnote{Condition que la page laisse
implicite et que cette lecture ajoute : sans elle, b) ne définirait pas une
propriété des morphismes de faisceaux indépendante de la présentation.}.

\pagerange{11}{11}
\section*{Fragment d'une démonstration (page 11)}

La page reprend au milieu d'une phrase ; le théorème démontré et le début de
la démonstration ne sont pas dans le dossier\footnote{La page est parmi les plus
difficiles du dossier : de nombreux mots y sont illisibles, des insertions
interlinéaires et des biffures se chevauchent au milieu, et l'ordre des mots
rendus n'y est pas assuré. Un bloc encadré, où se lisent « $Y$ intègre »,
« extension radicielle finie $Y'$ de $Y$ » et « normalisant », est biffé.
Ce qui suit reconstitue le mécanisme ; les liaisons logiques sont de cette
lecture.}. Ce qu'on en lit suffit à reconnaître une technique : rendre réduites
les fibres d'une famille, après une extension radicielle, à l'aide du schéma de
Hilbert. On y suppose, sans que la page le dise, $X \to Y$ assez régulier pour
que le schéma de Hilbert $\underline{\mathrm{Hilb}}_{X/Y}$ existe (par exemple
projectif et de présentation finie)\footnote{Hypothèse ajoutée.}.

\emph{Première étape.} Pour un point $y$ de $Y$, il existe une extension
radicielle finie $k'$ de $k(y)$ telle que $((X_y)_{k'})_{\mathrm{red}}$ soit
géométriquement réduit sur $k'$ (« séparable »\footnote{Lecture incertaine,
ainsi que « $k(y)$ ».}) ; quitte à faire ce changement de base et à remplacer $X$
par $X_{\mathrm{red}}$, on se ramène au cas où $(X_y)_{\mathrm{red}}$ est
géométriquement réduit, les hypothèses du théorème étant conservées.

\emph{Deuxième étape.} Une deuxième assertion du théorème se déduit d'une
troisième par réduction au cas noethérien. Pour celle-ci, on suppose $Y$ réduit.
Soit $M = \underline{\mathrm{Hilb}}_{X/Y}$, et $\varphi$ la section
\emph{ensembliste} de $M$ au-dessus de $Y$ qui envoie $y$ sur le point de la
fibre $M_y$ défini par le sous-schéma fermé $(X_y)_{\mathrm{red}} \subset X_y$.
Soit $Y'$ le sous-ensemble $\varphi(Y)$ de $M$. La page affirme, en
l'admettant, que $Y'$ est fermé dans $M$ et que $Y' \to Y$ est fini ; muni de sa
structure réduite, $Y'$ est alors un $Y$-schéma, et l'inclusion $Y' \to M$
correspond à un sous-schéma fermé $Z$ de $X' = X \times_Y Y'$, plat et de
présentation finie sur $Y'$\footnote{La page écrit « sur $Y$ » ; le sens demande
$Y'$.}, dont la fibre en $y'$ au-dessus de $y$ est
$(X_y)_{\mathrm{red}} \otimes_{k(y)} k(y')$. Ces fibres sont réduites — c'est
pour cela que la première étape rendait $(X_y)_{\mathrm{red}}$ géométriquement
réduit. Un morphisme plat, à fibres réduites, sur une base réduite a une source
réduite (EGA IV, dans le cas localement noethérien, cité de mémoire) ; donc $Z$ est réduit, et comme il a même
ensemble sous-jacent que $X'$, il est égal à $X'_{\mathrm{red}}$.

Le résultat est que, sur $Y'$, la réduction de $X'$ est plate à fibres
géométriquement réduites. C'est un énoncé du type « aplatissement de la
réduction après un changement de base fini radiciel » ; le dossier ne dit pas
dans quel théorème il s'insère.

\pagerange{12}{13}
\section*{Un critère d'immersion ouverte (pages 12 et 13)}

Soit $S$ un schéma localement noethérien et

\begin{tikzcd}
G \arrow[r, "u"] & F \arrow[d] \\
 & S
\end{tikzcd}

un morphisme de foncteurs $(\mathrm{Sch})^{\circ}_{/S} \to (\mathrm{Ens})$.
Supposons :
\begin{itemize}
  \item[(i)] $F$ et $G$ sont des faisceaux pour la topologie fpqc, commutent aux
    limites inductives filtrantes d'anneaux et aux limites projectives adiques
    d'anneaux artiniens ;
  \item[(ii)] $u$ est un monomorphisme, et formellement étale\footnote{La page
    dit « infinitésimalement étale ». Pour un monomorphisme, la condition de
    non-ramification est automatique, et il s'agit du relèvement le long des
    épaississements nilpotents. Une troisième hypothèse, « $u$ … épimorphisme »,
    est biffée ; « Si » est biffé et remplacé par « Supposons ».}.
\end{itemize}
Alors $u$ est une immersion ouverte : pour tout schéma $T$ et tout $T \to F$,
le sous-foncteur $G_T = G \times_F T$ de $T$ est représentable par un ouvert de
$T$.

La page énonce le critère sans démonstration. On peut le justifier
ainsi\footnote{Cette démonstration est de cette lecture.}. Par la commutation aux
limites inductives filtrantes, on se ramène à $T$ de type fini sur $S$, donc
localement noethérien. Soit $t \in T$ dont le point $\mathrm{Spec}\,k(t) \to T$
se factorise par $G_T$. Comme $u$ est formellement étale, ce point se relève de
façon unique aux $\mathrm{Spec}\,\mathcal{O}_{T,t}/\mathfrak{m}^n$, qui sont
artiniens ; par la commutation aux limites adiques, il se relève à
$\mathrm{Spec}\,\widehat{\mathcal{O}}_{T,t}$ ; comme $G$ est un faisceau fpqc et
$u$ un monomorphisme, il descend à $\mathrm{Spec}\,\mathcal{O}_{T,t}$, et, $G$
étant lfp, à un voisinage ouvert de $t$. La réunion $U$ des ouverts de $T$ qui se
factorisent par $G_T$ est donc un ouvert contenant tous ces points, et contenu
dans $G_T$ puisque $G_T$ est un faisceau ; inversement, tout $T' \to G_T$ envoie
les points de $T'$ dans $U$ (encore par descente fpqc le long des extensions
de corps), donc se factorise par l'ouvert $U$. D'où $G_T = U$.

Les deux énoncés qui suivent, séparés par des traits : une immersion ouverte
surjective de foncteurs est un isomorphisme (il suffit de le voir après
changement de base à un schéma, où c'est clair) ; et, annoncée seulement, la
caractérisation des morphismes étales de foncteurs.

Les conditions (i) sont celles qu'Artin a mises à la base de ses critères de
représentabilité — foncteur « limité » et effectivité des déformations
formelles\footnote{Nom moderne ; la page ne cite pas Artin ici. Les conditions
sont aussi celles que vérifie la page 4 du dossier 91 pour un théorème de
représentabilité « par morphismes non nécessairement séparés ».}. La marge de
la page 8 les annonçait déjà.

\textbf{La feuille de mots-clés.} La page 13 dresse, en cadres non reliés, la
carte des notions à traiter\footnote{Les cadres ne portent pas de flèches entre
eux, sauf une accolade ; la transcription les donne dans l'ordre de lecture.
« formellement » surcharge un « affine » biffé ; « immersion ouverte » est
biffé dans la liste des immersions et repris à côté ; « relatif » est une
lecture incertaine.} : l'immersion ouverte $X \to Y$ ; les morphismes
formellement lisses, non ramifiés, étales, et leurs versions relatives ; les
monomorphismes (« Mon ») et les immersions, fermées et ouvertes ; puis les
rubriques séparé, propre, connexe du plan de la page 10. Le carré de changement
de base

\begin{tikzcd}
G \arrow[r] & F \\
G_T \arrow[r] \arrow[u] & F_T \arrow[u]
\end{tikzcd}

rappelle le principe de vérification de la page 10, et une famille
$F_i \to F$ d'immersions ouvertes le recollement des pages 2 à 6. Enfin un
groupe de trois cadres réunis par une accolade et marqué « (A) » —
« faisceaux », « loc.\ de prés.\ finie », « adiques » — est la liste des
hypothèses (i) de la page 12.

\end{document}
